\begin{table}
\centering
\begin{tabular}{l>{[}l<{]}r}
\Hline
\textsc{layer} &  \mc{1}{l}{\textsc{span}} & \mc{1}{l}{\textsc{count}} \\
\Hline
1	&	1--22 & 69 \\
2	&	2--22 & 69 \\
3	&	3--21 & 69 \\
4	&	5--21 & 69 \\
5	&	5--19 & 69 \\
6	&	5--18 & 51 \\
22	&	13--15 & 40 \\
18	&	10--16 & 38 \\
7	&	5--17 & 37 \\
17	&	9--16 & 36 \\
25	&	9--10 & 27 \\
13	&	9--17 & 25 \\
24	&	6--8 & 24 \\
26	&	5--6 & 24 \\
8	&	6--17 & 23 \\
20	&	10--13 & 23 \\
14	&	8--16 & 20 \\
9	&	6--16 & 18 \\
11	&	8--17 & 18 \\
21	&	17--19 & 18 \\
10	&	9--18 & 14 \\
16	&	10--17 & 14 \\
23	&	8--10 & 11 \\
15	&	5--13 & 10 \\
19	&	6--10 & 8 \\
12	&	10--18 & 4 \\
\Hline
\end{tabular}
\caption{Number of trees each domain appears in within the constituency forest \label{ForestSpans}}
\end{table}
